% Options for packages loaded elsewhere
\PassOptionsToPackage{unicode}{hyperref}
\PassOptionsToPackage{hyphens}{url}
\documentclass[
]{article}
\usepackage{xcolor}
\usepackage[margin=1in]{geometry}
\usepackage{amsmath,amssymb}
\setcounter{secnumdepth}{-\maxdimen} % remove section numbering
\usepackage{iftex}
\ifPDFTeX
  \usepackage[T1]{fontenc}
  \usepackage[utf8]{inputenc}
  \usepackage{textcomp} % provide euro and other symbols
\else % if luatex or xetex
  \usepackage{unicode-math} % this also loads fontspec
  \defaultfontfeatures{Scale=MatchLowercase}
  \defaultfontfeatures[\rmfamily]{Ligatures=TeX,Scale=1}
\fi
\usepackage{lmodern}
\ifPDFTeX\else
  % xetex/luatex font selection
\fi
% Use upquote if available, for straight quotes in verbatim environments
\IfFileExists{upquote.sty}{\usepackage{upquote}}{}
\IfFileExists{microtype.sty}{% use microtype if available
  \usepackage[]{microtype}
  \UseMicrotypeSet[protrusion]{basicmath} % disable protrusion for tt fonts
}{}
\makeatletter
\@ifundefined{KOMAClassName}{% if non-KOMA class
  \IfFileExists{parskip.sty}{%
    \usepackage{parskip}
  }{% else
    \setlength{\parindent}{0pt}
    \setlength{\parskip}{6pt plus 2pt minus 1pt}}
}{% if KOMA class
  \KOMAoptions{parskip=half}}
\makeatother
\usepackage{color}
\usepackage{fancyvrb}
\newcommand{\VerbBar}{|}
\newcommand{\VERB}{\Verb[commandchars=\\\{\}]}
\DefineVerbatimEnvironment{Highlighting}{Verbatim}{commandchars=\\\{\}}
% Add ',fontsize=\small' for more characters per line
\usepackage{framed}
\definecolor{shadecolor}{RGB}{248,248,248}
\newenvironment{Shaded}{\begin{snugshade}}{\end{snugshade}}
\newcommand{\AlertTok}[1]{\textcolor[rgb]{0.94,0.16,0.16}{#1}}
\newcommand{\AnnotationTok}[1]{\textcolor[rgb]{0.56,0.35,0.01}{\textbf{\textit{#1}}}}
\newcommand{\AttributeTok}[1]{\textcolor[rgb]{0.13,0.29,0.53}{#1}}
\newcommand{\BaseNTok}[1]{\textcolor[rgb]{0.00,0.00,0.81}{#1}}
\newcommand{\BuiltInTok}[1]{#1}
\newcommand{\CharTok}[1]{\textcolor[rgb]{0.31,0.60,0.02}{#1}}
\newcommand{\CommentTok}[1]{\textcolor[rgb]{0.56,0.35,0.01}{\textit{#1}}}
\newcommand{\CommentVarTok}[1]{\textcolor[rgb]{0.56,0.35,0.01}{\textbf{\textit{#1}}}}
\newcommand{\ConstantTok}[1]{\textcolor[rgb]{0.56,0.35,0.01}{#1}}
\newcommand{\ControlFlowTok}[1]{\textcolor[rgb]{0.13,0.29,0.53}{\textbf{#1}}}
\newcommand{\DataTypeTok}[1]{\textcolor[rgb]{0.13,0.29,0.53}{#1}}
\newcommand{\DecValTok}[1]{\textcolor[rgb]{0.00,0.00,0.81}{#1}}
\newcommand{\DocumentationTok}[1]{\textcolor[rgb]{0.56,0.35,0.01}{\textbf{\textit{#1}}}}
\newcommand{\ErrorTok}[1]{\textcolor[rgb]{0.64,0.00,0.00}{\textbf{#1}}}
\newcommand{\ExtensionTok}[1]{#1}
\newcommand{\FloatTok}[1]{\textcolor[rgb]{0.00,0.00,0.81}{#1}}
\newcommand{\FunctionTok}[1]{\textcolor[rgb]{0.13,0.29,0.53}{\textbf{#1}}}
\newcommand{\ImportTok}[1]{#1}
\newcommand{\InformationTok}[1]{\textcolor[rgb]{0.56,0.35,0.01}{\textbf{\textit{#1}}}}
\newcommand{\KeywordTok}[1]{\textcolor[rgb]{0.13,0.29,0.53}{\textbf{#1}}}
\newcommand{\NormalTok}[1]{#1}
\newcommand{\OperatorTok}[1]{\textcolor[rgb]{0.81,0.36,0.00}{\textbf{#1}}}
\newcommand{\OtherTok}[1]{\textcolor[rgb]{0.56,0.35,0.01}{#1}}
\newcommand{\PreprocessorTok}[1]{\textcolor[rgb]{0.56,0.35,0.01}{\textit{#1}}}
\newcommand{\RegionMarkerTok}[1]{#1}
\newcommand{\SpecialCharTok}[1]{\textcolor[rgb]{0.81,0.36,0.00}{\textbf{#1}}}
\newcommand{\SpecialStringTok}[1]{\textcolor[rgb]{0.31,0.60,0.02}{#1}}
\newcommand{\StringTok}[1]{\textcolor[rgb]{0.31,0.60,0.02}{#1}}
\newcommand{\VariableTok}[1]{\textcolor[rgb]{0.00,0.00,0.00}{#1}}
\newcommand{\VerbatimStringTok}[1]{\textcolor[rgb]{0.31,0.60,0.02}{#1}}
\newcommand{\WarningTok}[1]{\textcolor[rgb]{0.56,0.35,0.01}{\textbf{\textit{#1}}}}
\usepackage{graphicx}
\makeatletter
\newsavebox\pandoc@box
\newcommand*\pandocbounded[1]{% scales image to fit in text height/width
  \sbox\pandoc@box{#1}%
  \Gscale@div\@tempa{\textheight}{\dimexpr\ht\pandoc@box+\dp\pandoc@box\relax}%
  \Gscale@div\@tempb{\linewidth}{\wd\pandoc@box}%
  \ifdim\@tempb\p@<\@tempa\p@\let\@tempa\@tempb\fi% select the smaller of both
  \ifdim\@tempa\p@<\p@\scalebox{\@tempa}{\usebox\pandoc@box}%
  \else\usebox{\pandoc@box}%
  \fi%
}
% Set default figure placement to htbp
\def\fps@figure{htbp}
\makeatother
\setlength{\emergencystretch}{3em} % prevent overfull lines
\providecommand{\tightlist}{%
  \setlength{\itemsep}{0pt}\setlength{\parskip}{0pt}}
\usepackage{graphicx}
\usepackage{float}
\usepackage{bookmark}
\IfFileExists{xurl.sty}{\usepackage{xurl}}{} % add URL line breaks if available
\urlstyle{same}
\hypersetup{
  hidelinks,
  pdfcreator={LaTeX via pandoc}}

\author{}
\date{\vspace{-2.5em}}

\begin{document}

\begin{Shaded}
\begin{Highlighting}[]
\FunctionTok{options}\NormalTok{(}\AttributeTok{OutDec =} \StringTok{"."}\NormalTok{)  }\CommentTok{\# Asegurar punto decimal en este chunk}

\CommentTok{\# Establecer semilla aleatoria}
\FunctionTok{set.seed}\NormalTok{(}\FunctionTok{sample}\NormalTok{(}\DecValTok{1}\SpecialCharTok{:}\DecValTok{10000}\NormalTok{, }\DecValTok{1}\NormalTok{))}

\CommentTok{\# Aleatorización del contexto del problema}
\NormalTok{contextos }\OtherTok{\textless{}{-}} \FunctionTok{c}\NormalTok{(}
  \StringTok{"empresa"}\NormalTok{, }\StringTok{"organización"}\NormalTok{, }\StringTok{"compañía"}\NormalTok{, }\StringTok{"institución"}\NormalTok{, }\StringTok{"corporación"}\NormalTok{,}
  \StringTok{"entidad"}\NormalTok{, }\StringTok{"firma"}\NormalTok{, }\StringTok{"negocio"}\NormalTok{, }\StringTok{"sociedad"}\NormalTok{, }\StringTok{"grupo empresarial"}
\NormalTok{)}
\NormalTok{contexto }\OtherTok{\textless{}{-}} \FunctionTok{sample}\NormalTok{(contextos, }\DecValTok{1}\NormalTok{)}

\CommentTok{\# Aleatorización de los tipos de bienes}
\NormalTok{tipos\_bien1 }\OtherTok{\textless{}{-}} \FunctionTok{c}\NormalTok{(}\StringTok{"carro"}\NormalTok{, }\StringTok{"auto"}\NormalTok{, }\StringTok{"vehículo"}\NormalTok{, }\StringTok{"automóvil"}\NormalTok{, }\StringTok{"coche"}\NormalTok{)}
\NormalTok{tipos\_bien2 }\OtherTok{\textless{}{-}} \FunctionTok{c}\NormalTok{(}\StringTok{"casa"}\NormalTok{, }\StringTok{"vivienda"}\NormalTok{, }\StringTok{"residencia"}\NormalTok{, }\StringTok{"hogar"}\NormalTok{, }\StringTok{"domicilio"}\NormalTok{)}
\NormalTok{tipos\_bien3 }\OtherTok{\textless{}{-}} \FunctionTok{c}\NormalTok{(}\StringTok{"apartamento"}\NormalTok{, }\StringTok{"departamento"}\NormalTok{, }\StringTok{"piso"}\NormalTok{, }\StringTok{"condominio"}\NormalTok{)}

\NormalTok{bien1 }\OtherTok{\textless{}{-}} \FunctionTok{sample}\NormalTok{(tipos\_bien1, }\DecValTok{1}\NormalTok{)}
\NormalTok{bien2 }\OtherTok{\textless{}{-}} \FunctionTok{sample}\NormalTok{(tipos\_bien2, }\DecValTok{1}\NormalTok{)}
\NormalTok{bien3 }\OtherTok{\textless{}{-}} \FunctionTok{sample}\NormalTok{(tipos\_bien3, }\DecValTok{1}\NormalTok{)}

\CommentTok{\# Aleatorización de términos para el enunciado}
\NormalTok{terminos\_encuesta }\OtherTok{\textless{}{-}} \FunctionTok{c}\NormalTok{(}\StringTok{"encuesta"}\NormalTok{, }\StringTok{"sondeo"}\NormalTok{, }\StringTok{"estudio"}\NormalTok{, }\StringTok{"investigación"}\NormalTok{, }\StringTok{"consulta"}\NormalTok{)}
\NormalTok{termino\_encuesta }\OtherTok{\textless{}{-}} \FunctionTok{sample}\NormalTok{(terminos\_encuesta, }\DecValTok{1}\NormalTok{)}

\NormalTok{terminos\_personas }\OtherTok{\textless{}{-}} \FunctionTok{c}\NormalTok{(}\StringTok{"personas"}\NormalTok{, }\StringTok{"empleados"}\NormalTok{, }\StringTok{"trabajadores"}\NormalTok{, }\StringTok{"colaboradores"}\NormalTok{, }\StringTok{"miembros"}\NormalTok{)}
\NormalTok{termino\_personas }\OtherTok{\textless{}{-}} \FunctionTok{sample}\NormalTok{(terminos\_personas, }\DecValTok{1}\NormalTok{)}

\NormalTok{terminos\_bienes }\OtherTok{\textless{}{-}} \FunctionTok{c}\NormalTok{(}\StringTok{"bienes"}\NormalTok{, }\StringTok{"posesiones"}\NormalTok{, }\StringTok{"propiedades"}\NormalTok{, }\StringTok{"activos"}\NormalTok{, }\StringTok{"patrimonios"}\NormalTok{)}
\NormalTok{termino\_bienes }\OtherTok{\textless{}{-}} \FunctionTok{sample}\NormalTok{(terminos\_bienes, }\DecValTok{1}\NormalTok{)}

\NormalTok{terminos\_resultados }\OtherTok{\textless{}{-}} \FunctionTok{c}\NormalTok{(}\StringTok{"resultados"}\NormalTok{, }\StringTok{"datos"}\NormalTok{, }\StringTok{"estadísticas"}\NormalTok{, }\StringTok{"cifras"}\NormalTok{)}
\NormalTok{termino\_resultados }\OtherTok{\textless{}{-}} \FunctionTok{sample}\NormalTok{(terminos\_resultados, }\DecValTok{1}\NormalTok{)}

\CommentTok{\# Aleatorización de colores para el gráfico circular (paletas oscuras)}
\NormalTok{paleta\_colores }\OtherTok{\textless{}{-}} \FunctionTok{list}\NormalTok{(}
  \FunctionTok{c}\NormalTok{(}\StringTok{"\#C62828"}\NormalTok{, }\StringTok{"\#2E7D32"}\NormalTok{, }\StringTok{"\#1565C0"}\NormalTok{, }\StringTok{"\#EF6C00"}\NormalTok{, }\StringTok{"\#6A1B9A"}\NormalTok{),  }\CommentTok{\# Paleta oscura 1}
  \FunctionTok{c}\NormalTok{(}\StringTok{"\#D32F2F"}\NormalTok{, }\StringTok{"\#388E3C"}\NormalTok{, }\StringTok{"\#1976D2"}\NormalTok{, }\StringTok{"\#F57C00"}\NormalTok{, }\StringTok{"\#7B1FA2"}\NormalTok{),  }\CommentTok{\# Paleta saturada}
  \FunctionTok{c}\NormalTok{(}\StringTok{"\#B71C1C"}\NormalTok{, }\StringTok{"\#1B5E20"}\NormalTok{, }\StringTok{"\#0D47A1"}\NormalTok{, }\StringTok{"\#E65100"}\NormalTok{, }\StringTok{"\#4A148C"}\NormalTok{),  }\CommentTok{\# Paleta oscura 2}
  \FunctionTok{c}\NormalTok{(}\StringTok{"\#AD1457"}\NormalTok{, }\StringTok{"\#00695C"}\NormalTok{, }\StringTok{"\#283593"}\NormalTok{, }\StringTok{"\#FF6F00"}\NormalTok{, }\StringTok{"\#4A148C"}\NormalTok{),  }\CommentTok{\# Paleta oscura 3}
  \FunctionTok{c}\NormalTok{(}\StringTok{"\#880E4F"}\NormalTok{, }\StringTok{"\#1B5E20"}\NormalTok{, }\StringTok{"\#0D47A1"}\NormalTok{, }\StringTok{"\#BF360C"}\NormalTok{, }\StringTok{"\#4A148C"}\NormalTok{)   }\CommentTok{\# Paleta oscura 4}
\NormalTok{)}
\NormalTok{paleta\_seleccionada }\OtherTok{\textless{}{-}} \FunctionTok{sample}\NormalTok{(paleta\_colores, }\DecValTok{1}\NormalTok{)[[}\DecValTok{1}\NormalTok{]]}

\CommentTok{\# Aleatorización de los porcentajes manteniendo coherencia matemática}
\CommentTok{\# Generación de porcentajes iniciales para cada categoría}
\NormalTok{set\_porcentajes }\OtherTok{\textless{}{-}} \ControlFlowTok{function}\NormalTok{() \{}
  \ControlFlowTok{repeat}\NormalTok{ \{}
    \CommentTok{\# Generar valores base aleatorios}
\NormalTok{    p\_bien1\_bien3 }\OtherTok{\textless{}{-}} \FunctionTok{sample}\NormalTok{(}\DecValTok{15}\SpecialCharTok{:}\DecValTok{35}\NormalTok{, }\DecValTok{1}\NormalTok{)  }\CommentTok{\# Porcentaje de bien1 y bien3}
\NormalTok{    p\_solo\_bien3 }\OtherTok{\textless{}{-}} \FunctionTok{sample}\NormalTok{(}\DecValTok{15}\SpecialCharTok{:}\DecValTok{30}\NormalTok{, }\DecValTok{1}\NormalTok{)   }\CommentTok{\# Porcentaje solo bien3}
\NormalTok{    p\_solo\_bien2 }\OtherTok{\textless{}{-}} \FunctionTok{sample}\NormalTok{(}\DecValTok{25}\SpecialCharTok{:}\DecValTok{40}\NormalTok{, }\DecValTok{1}\NormalTok{)   }\CommentTok{\# Porcentaje solo bien2}
\NormalTok{    p\_solo\_bien1 }\OtherTok{\textless{}{-}} \FunctionTok{sample}\NormalTok{(}\DecValTok{5}\SpecialCharTok{:}\DecValTok{10}\NormalTok{, }\DecValTok{1}\NormalTok{)    }\CommentTok{\# Porcentaje solo bien1}

    \CommentTok{\# Calcular el porcentaje restante para bien1 y bien2}
\NormalTok{    p\_bien1\_bien2 }\OtherTok{\textless{}{-}} \DecValTok{100} \SpecialCharTok{{-}}\NormalTok{ (p\_bien1\_bien3 }\SpecialCharTok{+}\NormalTok{ p\_solo\_bien3 }\SpecialCharTok{+}\NormalTok{ p\_solo\_bien2 }\SpecialCharTok{+}\NormalTok{ p\_solo\_bien1)}

    \CommentTok{\# Validar que el porcentaje restante sea positivo y razonable}
    \ControlFlowTok{if}\NormalTok{ (p\_bien1\_bien2 }\SpecialCharTok{\textgreater{}=} \DecValTok{10} \SpecialCharTok{\&\&}\NormalTok{ p\_bien1\_bien2 }\SpecialCharTok{\textless{}=} \DecValTok{25}\NormalTok{) \{}
      \FunctionTok{return}\NormalTok{(}\FunctionTok{c}\NormalTok{(p\_bien1\_bien3, p\_solo\_bien3, p\_solo\_bien2, p\_solo\_bien1, p\_bien1\_bien2))}
\NormalTok{    \}}
\NormalTok{  \}}
\NormalTok{\}}

\CommentTok{\# Obtener porcentajes válidos}
\NormalTok{porcentajes\_base }\OtherTok{\textless{}{-}} \FunctionTok{set\_porcentajes}\NormalTok{()}

\NormalTok{p\_bien1\_bien3 }\OtherTok{\textless{}{-}}\NormalTok{ porcentajes\_base[}\DecValTok{1}\NormalTok{]  }\CommentTok{\# Porcentaje de bien1 y bien3}
\NormalTok{p\_solo\_bien3 }\OtherTok{\textless{}{-}}\NormalTok{ porcentajes\_base[}\DecValTok{2}\NormalTok{]   }\CommentTok{\# Porcentaje solo bien3}
\NormalTok{p\_solo\_bien2 }\OtherTok{\textless{}{-}}\NormalTok{ porcentajes\_base[}\DecValTok{3}\NormalTok{]   }\CommentTok{\# Porcentaje solo bien2}
\NormalTok{p\_solo\_bien1 }\OtherTok{\textless{}{-}}\NormalTok{ porcentajes\_base[}\DecValTok{4}\NormalTok{]   }\CommentTok{\# Porcentaje solo bien1}
\NormalTok{p\_bien1\_bien2 }\OtherTok{\textless{}{-}}\NormalTok{ porcentajes\_base[}\DecValTok{5}\NormalTok{]  }\CommentTok{\# Porcentaje de bien1 y bien2}

\CommentTok{\# Guardar los porcentajes originales para visualización}
\NormalTok{porcentajes\_grafico }\OtherTok{\textless{}{-}} \FunctionTok{c}\NormalTok{(p\_bien1\_bien3, p\_solo\_bien3, p\_solo\_bien2, p\_solo\_bien1, p\_bien1\_bien2)}

\CommentTok{\# Verificar que suman 100\%}
\FunctionTok{test\_that}\NormalTok{(}\StringTok{"Los porcentajes suman 100\%"}\NormalTok{, \{}
  \FunctionTok{expect\_equal}\NormalTok{(}\FunctionTok{sum}\NormalTok{(porcentajes\_base), }\DecValTok{100}\NormalTok{)}
\NormalTok{\})}
\end{Highlighting}
\end{Shaded}

Test passed

\begin{Shaded}
\begin{Highlighting}[]
\CommentTok{\# ENFOQUE MEJORADO: Primero determinamos el total y luego las categorías sin redondear}
\CommentTok{\# Número total de personas {-} elegimos un número que sea divisible por 100 para evitar errores de redondeo}
\NormalTok{total\_personas }\OtherTok{\textless{}{-}} \FunctionTok{sample}\NormalTok{(}\FunctionTok{c}\NormalTok{(}\DecValTok{500}\NormalTok{, }\DecValTok{600}\NormalTok{, }\DecValTok{700}\NormalTok{, }\DecValTok{800}\NormalTok{, }\DecValTok{900}\NormalTok{, }\DecValTok{1000}\NormalTok{, }\DecValTok{1200}\NormalTok{, }\DecValTok{1500}\NormalTok{, }\DecValTok{2000}\NormalTok{), }\DecValTok{1}\NormalTok{)}

\CommentTok{\# Calculamos la cantidad exacta de personas para cada categoría (sin redondeo)}
\CommentTok{\# Usamos división por enteros para evitar problemas decimales}
\NormalTok{personas\_exactas }\OtherTok{\textless{}{-}} \FunctionTok{list}\NormalTok{(}
  \AttributeTok{bien1\_bien3 =}\NormalTok{ (total\_personas }\SpecialCharTok{*}\NormalTok{ p\_bien1\_bien3) }\SpecialCharTok{\%/\%} \DecValTok{100}\NormalTok{,}
  \AttributeTok{solo\_bien3 =}\NormalTok{ (total\_personas }\SpecialCharTok{*}\NormalTok{ p\_solo\_bien3) }\SpecialCharTok{\%/\%} \DecValTok{100}\NormalTok{,}
  \AttributeTok{solo\_bien2 =}\NormalTok{ (total\_personas }\SpecialCharTok{*}\NormalTok{ p\_solo\_bien2) }\SpecialCharTok{\%/\%} \DecValTok{100}\NormalTok{,}
  \AttributeTok{solo\_bien1 =}\NormalTok{ (total\_personas }\SpecialCharTok{*}\NormalTok{ p\_solo\_bien1) }\SpecialCharTok{\%/\%} \DecValTok{100}\NormalTok{,}
  \AttributeTok{bien1\_bien2 =}\NormalTok{ (total\_personas }\SpecialCharTok{*}\NormalTok{ p\_bien1\_bien2) }\SpecialCharTok{\%/\%} \DecValTok{100}
\NormalTok{)}

\CommentTok{\# Calculamos el residuo (personas no asignadas por la división entera)}
\NormalTok{personas\_no\_asignadas }\OtherTok{\textless{}{-}}\NormalTok{ total\_personas }\SpecialCharTok{{-}} \FunctionTok{sum}\NormalTok{(}\FunctionTok{unlist}\NormalTok{(personas\_exactas))}

\CommentTok{\# Distribuimos las personas no asignadas a las categorías en orden}
\CommentTok{\# hasta completar exactamente el total}
\ControlFlowTok{if}\NormalTok{ (personas\_no\_asignadas }\SpecialCharTok{\textgreater{}} \DecValTok{0}\NormalTok{) \{}
\NormalTok{  categorias }\OtherTok{\textless{}{-}} \FunctionTok{c}\NormalTok{(}\StringTok{"bien1\_bien3"}\NormalTok{, }\StringTok{"solo\_bien3"}\NormalTok{, }\StringTok{"solo\_bien2"}\NormalTok{, }\StringTok{"solo\_bien1"}\NormalTok{, }\StringTok{"bien1\_bien2"}\NormalTok{)}
  \ControlFlowTok{for}\NormalTok{ (i }\ControlFlowTok{in} \DecValTok{1}\SpecialCharTok{:}\NormalTok{personas\_no\_asignadas) \{}
\NormalTok{    idx }\OtherTok{\textless{}{-}}\NormalTok{ ((i}\DecValTok{{-}1}\NormalTok{) }\SpecialCharTok{\%\%} \DecValTok{5}\NormalTok{) }\SpecialCharTok{+} \DecValTok{1}  \CommentTok{\# Cicla entre 1 y 5}
\NormalTok{    personas\_exactas[[categorias[idx]]] }\OtherTok{\textless{}{-}}\NormalTok{ personas\_exactas[[categorias[idx]]] }\SpecialCharTok{+} \DecValTok{1}
\NormalTok{  \}}
\NormalTok{\}}

\CommentTok{\# Asignamos los valores finales de personas por categoría}
\NormalTok{personas\_bien1\_bien3 }\OtherTok{\textless{}{-}}\NormalTok{ personas\_exactas}\SpecialCharTok{$}\NormalTok{bien1\_bien3}
\NormalTok{personas\_solo\_bien3 }\OtherTok{\textless{}{-}}\NormalTok{ personas\_exactas}\SpecialCharTok{$}\NormalTok{solo\_bien3}
\NormalTok{personas\_solo\_bien2 }\OtherTok{\textless{}{-}}\NormalTok{ personas\_exactas}\SpecialCharTok{$}\NormalTok{solo\_bien2}
\NormalTok{personas\_solo\_bien1 }\OtherTok{\textless{}{-}}\NormalTok{ personas\_exactas}\SpecialCharTok{$}\NormalTok{solo\_bien1}
\NormalTok{personas\_bien1\_bien2 }\OtherTok{\textless{}{-}}\NormalTok{ personas\_exactas}\SpecialCharTok{$}\NormalTok{bien1\_bien2}

\CommentTok{\# Verificamos que la suma sea exactamente igual al total}
\NormalTok{personas }\OtherTok{\textless{}{-}} \FunctionTok{c}\NormalTok{(personas\_bien1\_bien3, personas\_solo\_bien3, personas\_solo\_bien2,}
\NormalTok{              personas\_solo\_bien1, personas\_bien1\_bien2)}
\FunctionTok{test\_that}\NormalTok{(}\StringTok{"La suma de todas las categorías es igual al total"}\NormalTok{, \{}
  \FunctionTok{expect\_equal}\NormalTok{(}\FunctionTok{sum}\NormalTok{(personas), total\_personas)}
\NormalTok{\})}
\end{Highlighting}
\end{Shaded}

Test passed

\begin{Shaded}
\begin{Highlighting}[]
\CommentTok{\# Recalculamos los porcentajes exactos basados en las cantidades finales de personas}
\CommentTok{\# Estos son los porcentajes que usaremos para todos los cálculos en la solución}
\NormalTok{porcentajes\_exactos }\OtherTok{\textless{}{-}} \FunctionTok{list}\NormalTok{(}
  \AttributeTok{bien1\_bien3 =}\NormalTok{ (personas\_bien1\_bien3 }\SpecialCharTok{*} \DecValTok{100}\NormalTok{) }\SpecialCharTok{/}\NormalTok{ total\_personas,}
  \AttributeTok{solo\_bien3 =}\NormalTok{ (personas\_solo\_bien3 }\SpecialCharTok{*} \DecValTok{100}\NormalTok{) }\SpecialCharTok{/}\NormalTok{ total\_personas,}
  \AttributeTok{solo\_bien2 =}\NormalTok{ (personas\_solo\_bien2 }\SpecialCharTok{*} \DecValTok{100}\NormalTok{) }\SpecialCharTok{/}\NormalTok{ total\_personas,}
  \AttributeTok{solo\_bien1 =}\NormalTok{ (personas\_solo\_bien1 }\SpecialCharTok{*} \DecValTok{100}\NormalTok{) }\SpecialCharTok{/}\NormalTok{ total\_personas,}
  \AttributeTok{bien1\_bien2 =}\NormalTok{ (personas\_bien1\_bien2 }\SpecialCharTok{*} \DecValTok{100}\NormalTok{) }\SpecialCharTok{/}\NormalTok{ total\_personas}
\NormalTok{)}

\CommentTok{\# Y los porcentajes redondeados para mostrar en pantalla}
\NormalTok{porcentajes\_mostrar }\OtherTok{\textless{}{-}} \FunctionTok{list}\NormalTok{(}
  \AttributeTok{bien1\_bien3 =} \FunctionTok{round}\NormalTok{(porcentajes\_exactos}\SpecialCharTok{$}\NormalTok{bien1\_bien3),}
  \AttributeTok{solo\_bien3 =} \FunctionTok{round}\NormalTok{(porcentajes\_exactos}\SpecialCharTok{$}\NormalTok{solo\_bien3),}
  \AttributeTok{solo\_bien2 =} \FunctionTok{round}\NormalTok{(porcentajes\_exactos}\SpecialCharTok{$}\NormalTok{solo\_bien2),}
  \AttributeTok{solo\_bien1 =} \FunctionTok{round}\NormalTok{(porcentajes\_exactos}\SpecialCharTok{$}\NormalTok{solo\_bien1),}
  \AttributeTok{bien1\_bien2 =} \FunctionTok{round}\NormalTok{(porcentajes\_exactos}\SpecialCharTok{$}\NormalTok{bien1\_bien2)}
\NormalTok{)}

\CommentTok{\# Verificamos que el total de los porcentajes exactos sea 100\%}
\FunctionTok{test\_that}\NormalTok{(}\StringTok{"Los porcentajes exactos suman 100\%"}\NormalTok{, \{}
  \FunctionTok{expect\_true}\NormalTok{(}\FunctionTok{abs}\NormalTok{(}\FunctionTok{sum}\NormalTok{(}\FunctionTok{unlist}\NormalTok{(porcentajes\_exactos)) }\SpecialCharTok{{-}} \DecValTok{100}\NormalTok{) }\SpecialCharTok{\textless{}} \FloatTok{0.001}\NormalTok{)}
\NormalTok{\})}
\end{Highlighting}
\end{Shaded}

Test passed

\begin{Shaded}
\begin{Highlighting}[]
\CommentTok{\# Para visualización y verificación}
\FunctionTok{cat}\NormalTok{(}\StringTok{"Total de personas:"}\NormalTok{, total\_personas, }\StringTok{"}\SpecialCharTok{\textbackslash{}n}\StringTok{"}\NormalTok{)}
\end{Highlighting}
\end{Shaded}

Total de personas: 700

\begin{Shaded}
\begin{Highlighting}[]
\FunctionTok{cat}\NormalTok{(}\StringTok{"Distribución por categoría:}\SpecialCharTok{\textbackslash{}n}\StringTok{"}\NormalTok{)}
\end{Highlighting}
\end{Shaded}

Distribución por categoría:

\begin{Shaded}
\begin{Highlighting}[]
\FunctionTok{cat}\NormalTok{(}\StringTok{"{-} bien1\_bien3:"}\NormalTok{, personas\_bien1\_bien3, }\StringTok{"personas,"}\NormalTok{,}
\NormalTok{    porcentajes\_exactos}\SpecialCharTok{$}\NormalTok{bien1\_bien3, }\StringTok{"\% (exacto),"}\NormalTok{,}
\NormalTok{    porcentajes\_mostrar}\SpecialCharTok{$}\NormalTok{bien1\_bien3, }\StringTok{"\% (mostrado)}\SpecialCharTok{\textbackslash{}n}\StringTok{"}\NormalTok{)}
\end{Highlighting}
\end{Shaded}

\begin{itemize}
\tightlist
\item
  bien1\_bien3: 161 personas, 23 \% (exacto), 23 \% (mostrado)
\end{itemize}

\begin{Shaded}
\begin{Highlighting}[]
\FunctionTok{cat}\NormalTok{(}\StringTok{"{-} solo\_bien3:"}\NormalTok{, personas\_solo\_bien3, }\StringTok{"personas,"}\NormalTok{,}
\NormalTok{    porcentajes\_exactos}\SpecialCharTok{$}\NormalTok{solo\_bien3, }\StringTok{"\% (exacto),"}\NormalTok{,}
\NormalTok{    porcentajes\_mostrar}\SpecialCharTok{$}\NormalTok{solo\_bien3, }\StringTok{"\% (mostrado)}\SpecialCharTok{\textbackslash{}n}\StringTok{"}\NormalTok{)}
\end{Highlighting}
\end{Shaded}

\begin{itemize}
\tightlist
\item
  solo\_bien3: 161 personas, 23 \% (exacto), 23 \% (mostrado)
\end{itemize}

\begin{Shaded}
\begin{Highlighting}[]
\FunctionTok{cat}\NormalTok{(}\StringTok{"{-} solo\_bien2:"}\NormalTok{, personas\_solo\_bien2, }\StringTok{"personas,"}\NormalTok{,}
\NormalTok{    porcentajes\_exactos}\SpecialCharTok{$}\NormalTok{solo\_bien2, }\StringTok{"\% (exacto),"}\NormalTok{,}
\NormalTok{    porcentajes\_mostrar}\SpecialCharTok{$}\NormalTok{solo\_bien2, }\StringTok{"\% (mostrado)}\SpecialCharTok{\textbackslash{}n}\StringTok{"}\NormalTok{)}
\end{Highlighting}
\end{Shaded}

\begin{itemize}
\tightlist
\item
  solo\_bien2: 210 personas, 30 \% (exacto), 30 \% (mostrado)
\end{itemize}

\begin{Shaded}
\begin{Highlighting}[]
\FunctionTok{cat}\NormalTok{(}\StringTok{"{-} solo\_bien1:"}\NormalTok{, personas\_solo\_bien1, }\StringTok{"personas,"}\NormalTok{,}
\NormalTok{    porcentajes\_exactos}\SpecialCharTok{$}\NormalTok{solo\_bien1, }\StringTok{"\% (exacto),"}\NormalTok{,}
\NormalTok{    porcentajes\_mostrar}\SpecialCharTok{$}\NormalTok{solo\_bien1, }\StringTok{"\% (mostrado)}\SpecialCharTok{\textbackslash{}n}\StringTok{"}\NormalTok{)}
\end{Highlighting}
\end{Shaded}

\begin{itemize}
\tightlist
\item
  solo\_bien1: 63 personas, 9 \% (exacto), 9 \% (mostrado)
\end{itemize}

\begin{Shaded}
\begin{Highlighting}[]
\FunctionTok{cat}\NormalTok{(}\StringTok{"{-} bien1\_bien2:"}\NormalTok{, personas\_bien1\_bien2, }\StringTok{"personas,"}\NormalTok{,}
\NormalTok{    porcentajes\_exactos}\SpecialCharTok{$}\NormalTok{bien1\_bien2, }\StringTok{"\% (exacto),"}\NormalTok{,}
\NormalTok{    porcentajes\_mostrar}\SpecialCharTok{$}\NormalTok{bien1\_bien2, }\StringTok{"\% (mostrado)}\SpecialCharTok{\textbackslash{}n}\StringTok{"}\NormalTok{)}
\end{Highlighting}
\end{Shaded}

\begin{itemize}
\tightlist
\item
  bien1\_bien2: 105 personas, 15 \% (exacto), 15 \% (mostrado)
\end{itemize}

\begin{Shaded}
\begin{Highlighting}[]
\CommentTok{\# Mapa de datos para acceso uniforme a todas las categorías}
\NormalTok{datos\_categorias }\OtherTok{\textless{}{-}} \FunctionTok{list}\NormalTok{(}
  \AttributeTok{solo\_bien1 =} \FunctionTok{list}\NormalTok{(}
    \AttributeTok{personas =}\NormalTok{ personas\_solo\_bien1,}
    \AttributeTok{texto =} \FunctionTok{paste0}\NormalTok{(}\StringTok{"solo "}\NormalTok{, bien1),}
    \AttributeTok{porcentaje\_exacto =}\NormalTok{ porcentajes\_exactos}\SpecialCharTok{$}\NormalTok{solo\_bien1,}
    \AttributeTok{porcentaje\_mostrar =}\NormalTok{ porcentajes\_mostrar}\SpecialCharTok{$}\NormalTok{solo\_bien1}
\NormalTok{  ),}
  \AttributeTok{solo\_bien2 =} \FunctionTok{list}\NormalTok{(}
    \AttributeTok{personas =}\NormalTok{ personas\_solo\_bien2,}
    \AttributeTok{texto =} \FunctionTok{paste0}\NormalTok{(}\StringTok{"solo "}\NormalTok{, bien2),}
    \AttributeTok{porcentaje\_exacto =}\NormalTok{ porcentajes\_exactos}\SpecialCharTok{$}\NormalTok{solo\_bien2,}
    \AttributeTok{porcentaje\_mostrar =}\NormalTok{ porcentajes\_mostrar}\SpecialCharTok{$}\NormalTok{solo\_bien2}
\NormalTok{  ),}
  \AttributeTok{solo\_bien3 =} \FunctionTok{list}\NormalTok{(}
    \AttributeTok{personas =}\NormalTok{ personas\_solo\_bien3,}
    \AttributeTok{texto =} \FunctionTok{paste0}\NormalTok{(}\StringTok{"solo "}\NormalTok{, bien3),}
    \AttributeTok{porcentaje\_exacto =}\NormalTok{ porcentajes\_exactos}\SpecialCharTok{$}\NormalTok{solo\_bien3,}
    \AttributeTok{porcentaje\_mostrar =}\NormalTok{ porcentajes\_mostrar}\SpecialCharTok{$}\NormalTok{solo\_bien3}
\NormalTok{  ),}
  \AttributeTok{bien1\_bien2 =} \FunctionTok{list}\NormalTok{(}
    \AttributeTok{personas =}\NormalTok{ personas\_bien1\_bien2,}
    \AttributeTok{texto =} \FunctionTok{paste0}\NormalTok{(bien1, }\StringTok{" y "}\NormalTok{, bien2),}
    \AttributeTok{porcentaje\_exacto =}\NormalTok{ porcentajes\_exactos}\SpecialCharTok{$}\NormalTok{bien1\_bien2,}
    \AttributeTok{porcentaje\_mostrar =}\NormalTok{ porcentajes\_mostrar}\SpecialCharTok{$}\NormalTok{bien1\_bien2}
\NormalTok{  ),}
  \AttributeTok{bien1\_bien3 =} \FunctionTok{list}\NormalTok{(}
    \AttributeTok{personas =}\NormalTok{ personas\_bien1\_bien3,}
    \AttributeTok{texto =} \FunctionTok{paste0}\NormalTok{(bien1, }\StringTok{" y "}\NormalTok{, bien3),}
    \AttributeTok{porcentaje\_exacto =}\NormalTok{ porcentajes\_exactos}\SpecialCharTok{$}\NormalTok{bien1\_bien3,}
    \AttributeTok{porcentaje\_mostrar =}\NormalTok{ porcentajes\_mostrar}\SpecialCharTok{$}\NormalTok{bien1\_bien3}
\NormalTok{  )}
\NormalTok{)}

\CommentTok{\# Seleccionar un valor conocido para el enunciado}
\NormalTok{tipos\_condicion }\OtherTok{\textless{}{-}} \FunctionTok{c}\NormalTok{(}\StringTok{"bien1\_bien3"}\NormalTok{, }\StringTok{"solo\_bien3"}\NormalTok{, }\StringTok{"solo\_bien2"}\NormalTok{, }\StringTok{"solo\_bien1"}\NormalTok{, }\StringTok{"bien1\_bien2"}\NormalTok{)}
\NormalTok{tipo\_condicion }\OtherTok{\textless{}{-}} \FunctionTok{sample}\NormalTok{(tipos\_condicion, }\DecValTok{1}\NormalTok{)}

\CommentTok{\# Determinar el valor y texto de la condición inicial}
\NormalTok{valor\_condicion }\OtherTok{\textless{}{-}}\NormalTok{ datos\_categorias[[tipo\_condicion]]}\SpecialCharTok{$}\NormalTok{personas}
\NormalTok{texto\_condicion }\OtherTok{\textless{}{-}}\NormalTok{ datos\_categorias[[tipo\_condicion]]}\SpecialCharTok{$}\NormalTok{texto}
\NormalTok{porcentaje\_condicion\_exacto }\OtherTok{\textless{}{-}}\NormalTok{ datos\_categorias[[tipo\_condicion]]}\SpecialCharTok{$}\NormalTok{porcentaje\_exacto}
\NormalTok{porcentaje\_condicion\_mostrar }\OtherTok{\textless{}{-}}\NormalTok{ datos\_categorias[[tipo\_condicion]]}\SpecialCharTok{$}\NormalTok{porcentaje\_mostrar}

\CommentTok{\# Aleatorización para las preguntas en formato cloze}
\CommentTok{\# Pregunta 1: Cuántas personas poseen un bien específico}
\CommentTok{\# Excluimos el tipo\_condicion para evitar preguntas redundantes}
\NormalTok{bienes\_pregunta1 }\OtherTok{\textless{}{-}} \FunctionTok{c}\NormalTok{(}\StringTok{"solo\_bien1"}\NormalTok{, }\StringTok{"solo\_bien2"}\NormalTok{, }\StringTok{"solo\_bien3"}\NormalTok{, }\StringTok{"bien1\_bien2"}\NormalTok{, }\StringTok{"bien1\_bien3"}\NormalTok{)}
\NormalTok{bienes\_pregunta1 }\OtherTok{\textless{}{-}}\NormalTok{ bienes\_pregunta1[bienes\_pregunta1 }\SpecialCharTok{!=}\NormalTok{ tipo\_condicion]}
\NormalTok{tipo\_pregunta1 }\OtherTok{\textless{}{-}} \FunctionTok{sample}\NormalTok{(bienes\_pregunta1, }\DecValTok{1}\NormalTok{)}

\CommentTok{\# Determinar el texto y la respuesta para la pregunta 1}
\NormalTok{respuesta1 }\OtherTok{\textless{}{-}}\NormalTok{ datos\_categorias[[tipo\_pregunta1]]}\SpecialCharTok{$}\NormalTok{personas}
\NormalTok{texto\_pregunta1 }\OtherTok{\textless{}{-}}\NormalTok{ datos\_categorias[[tipo\_pregunta1]]}\SpecialCharTok{$}\NormalTok{texto}
\NormalTok{porcentaje\_pregunta1\_exacto }\OtherTok{\textless{}{-}}\NormalTok{ datos\_categorias[[tipo\_pregunta1]]}\SpecialCharTok{$}\NormalTok{porcentaje\_exacto}
\NormalTok{porcentaje\_pregunta1\_mostrar }\OtherTok{\textless{}{-}}\NormalTok{ datos\_categorias[[tipo\_pregunta1]]}\SpecialCharTok{$}\NormalTok{porcentaje\_mostrar}

\CommentTok{\# Pregunta 2: Probabilidad de tener un bien específico}
\CommentTok{\# Excluimos tanto tipo\_condicion como tipo\_pregunta1 para evitar redundancias}
\NormalTok{bienes\_pregunta2 }\OtherTok{\textless{}{-}} \FunctionTok{c}\NormalTok{(}\StringTok{"solo\_bien1"}\NormalTok{, }\StringTok{"solo\_bien2"}\NormalTok{, }\StringTok{"solo\_bien3"}\NormalTok{, }\StringTok{"bien1\_bien2"}\NormalTok{, }\StringTok{"bien1\_bien3"}\NormalTok{)}
\NormalTok{bienes\_pregunta2 }\OtherTok{\textless{}{-}}\NormalTok{ bienes\_pregunta2[}\SpecialCharTok{!}\NormalTok{(bienes\_pregunta2 }\SpecialCharTok{\%in\%} \FunctionTok{c}\NormalTok{(tipo\_condicion, tipo\_pregunta1))]}
\NormalTok{tipo\_pregunta2 }\OtherTok{\textless{}{-}} \FunctionTok{sample}\NormalTok{(bienes\_pregunta2, }\DecValTok{1}\NormalTok{)}

\CommentTok{\# Determinar el texto y la respuesta para la pregunta 2}
\NormalTok{texto\_pregunta2 }\OtherTok{\textless{}{-}}\NormalTok{ datos\_categorias[[tipo\_pregunta2]]}\SpecialCharTok{$}\NormalTok{texto}
\CommentTok{\# Calculamos la probabilidad exacta como la proporción de personas}
\NormalTok{probabilidad\_exacta }\OtherTok{\textless{}{-}}\NormalTok{ datos\_categorias[[tipo\_pregunta2]]}\SpecialCharTok{$}\NormalTok{personas }\SpecialCharTok{/}\NormalTok{ total\_personas}
\CommentTok{\# Redondeamos a 4 cifras significativas para cálculos internos}
\NormalTok{respuesta2 }\OtherTok{\textless{}{-}} \FunctionTok{signif}\NormalTok{(probabilidad\_exacta, }\DecValTok{4}\NormalTok{)}

\CommentTok{\# Pregunta 3: Total de personas encuestadas}
\NormalTok{respuesta3 }\OtherTok{\textless{}{-}}\NormalTok{ total\_personas}
\end{Highlighting}
\end{Shaded}

\begin{Shaded}
\begin{Highlighting}[]
\FunctionTok{options}\NormalTok{(}\AttributeTok{OutDec =} \StringTok{"."}\NormalTok{)  }\CommentTok{\# Asegurar punto decimal en este chunk}

\CommentTok{\# Aleatorizar estilos visuales adicionales para el gráfico}
\NormalTok{estilo\_grafico }\OtherTok{\textless{}{-}} \FunctionTok{list}\NormalTok{(}
  \CommentTok{\# Aleatorizar el ángulo de inicio}
  \AttributeTok{angulo\_inicio =} \FunctionTok{sample}\NormalTok{(}\FunctionTok{c}\NormalTok{(}\DecValTok{0}\NormalTok{, }\DecValTok{45}\NormalTok{, }\DecValTok{90}\NormalTok{, }\DecValTok{135}\NormalTok{, }\DecValTok{180}\NormalTok{, }\DecValTok{225}\NormalTok{, }\DecValTok{270}\NormalTok{, }\DecValTok{315}\NormalTok{), }\DecValTok{1}\NormalTok{),}

  \CommentTok{\# Aleatorizar el estilo de borde}
  \AttributeTok{borde =} \FunctionTok{sample}\NormalTok{(}\FunctionTok{list}\NormalTok{(}
    \FunctionTok{list}\NormalTok{(}\AttributeTok{color =} \StringTok{"white"}\NormalTok{, }\AttributeTok{grosor =} \DecValTok{1}\NormalTok{),}
    \FunctionTok{list}\NormalTok{(}\AttributeTok{color =} \StringTok{"white"}\NormalTok{, }\AttributeTok{grosor =} \FloatTok{1.5}\NormalTok{),}
    \FunctionTok{list}\NormalTok{(}\AttributeTok{color =} \StringTok{"\#DDDDDD"}\NormalTok{, }\AttributeTok{grosor =} \DecValTok{1}\NormalTok{),}
    \FunctionTok{list}\NormalTok{(}\AttributeTok{color =} \StringTok{"\#333333"}\NormalTok{, }\AttributeTok{grosor =} \FloatTok{0.5}\NormalTok{),}
    \FunctionTok{list}\NormalTok{(}\AttributeTok{color =} \StringTok{"none"}\NormalTok{, }\AttributeTok{grosor =} \DecValTok{0}\NormalTok{)}
\NormalTok{  ), }\DecValTok{1}\NormalTok{)[[}\DecValTok{1}\NormalTok{]],}

  \CommentTok{\# Aleatorizar el estilo de sombra}
  \AttributeTok{sombra =} \FunctionTok{sample}\NormalTok{(}\FunctionTok{c}\NormalTok{(}\ConstantTok{TRUE}\NormalTok{, }\ConstantTok{FALSE}\NormalTok{), }\DecValTok{1}\NormalTok{, }\AttributeTok{prob =} \FunctionTok{c}\NormalTok{(}\FloatTok{0.7}\NormalTok{, }\FloatTok{0.3}\NormalTok{)),}

  \CommentTok{\# Aleatorizar el estilo de etiquetas de porcentaje}
  \AttributeTok{estilo\_porcentaje =} \FunctionTok{sample}\NormalTok{(}\FunctionTok{list}\NormalTok{(}
    \FunctionTok{list}\NormalTok{(}\AttributeTok{color\_fondo =} \StringTok{"black"}\NormalTok{, }\AttributeTok{alpha =} \FloatTok{0.7}\NormalTok{, }\AttributeTok{estilo =} \StringTok{"round"}\NormalTok{, }\AttributeTok{padding =} \FloatTok{0.2}\NormalTok{),}
    \FunctionTok{list}\NormalTok{(}\AttributeTok{color\_fondo =} \StringTok{"\#333333"}\NormalTok{, }\AttributeTok{alpha =} \FloatTok{0.8}\NormalTok{, }\AttributeTok{estilo =} \StringTok{"round"}\NormalTok{, }\AttributeTok{padding =} \FloatTok{0.3}\NormalTok{),}
    \FunctionTok{list}\NormalTok{(}\AttributeTok{color\_fondo =} \StringTok{"\#555555"}\NormalTok{, }\AttributeTok{alpha =} \FloatTok{0.7}\NormalTok{, }\AttributeTok{estilo =} \StringTok{"round,pad=0.2"}\NormalTok{, }\AttributeTok{padding =} \FloatTok{0.2}\NormalTok{),}
    \FunctionTok{list}\NormalTok{(}\AttributeTok{color\_fondo =} \StringTok{"none"}\NormalTok{, }\AttributeTok{alpha =} \DecValTok{1}\NormalTok{, }\AttributeTok{estilo =} \StringTok{"round"}\NormalTok{, }\AttributeTok{padding =} \DecValTok{0}\NormalTok{)}
\NormalTok{  ), }\DecValTok{1}\NormalTok{)[[}\DecValTok{1}\NormalTok{]],}

  \CommentTok{\# Aleatorizar el estilo de las etiquetas externas}
  \AttributeTok{estilo\_etiquetas =} \FunctionTok{sample}\NormalTok{(}\FunctionTok{list}\NormalTok{(}
    \FunctionTok{list}\NormalTok{(}\AttributeTok{color\_fondo =} \StringTok{"white"}\NormalTok{, }\AttributeTok{color\_borde =} \StringTok{"gray"}\NormalTok{, }\AttributeTok{alpha =} \FloatTok{0.9}\NormalTok{, }\AttributeTok{estilo =} \StringTok{"round"}\NormalTok{, }\AttributeTok{padding =} \FloatTok{0.3}\NormalTok{),}
    \FunctionTok{list}\NormalTok{(}\AttributeTok{color\_fondo =} \StringTok{"white"}\NormalTok{, }\AttributeTok{color\_borde =} \StringTok{"\#AAAAAA"}\NormalTok{, }\AttributeTok{alpha =} \FloatTok{0.95}\NormalTok{, }\AttributeTok{estilo =} \StringTok{"round"}\NormalTok{, }\AttributeTok{padding =} \FloatTok{0.4}\NormalTok{),}
    \FunctionTok{list}\NormalTok{(}\AttributeTok{color\_fondo =} \StringTok{"\#F5F5F5"}\NormalTok{, }\AttributeTok{color\_borde =} \StringTok{"\#CCCCCC"}\NormalTok{, }\AttributeTok{alpha =} \FloatTok{0.9}\NormalTok{, }\AttributeTok{estilo =} \StringTok{"round"}\NormalTok{, }\AttributeTok{padding =} \FloatTok{0.3}\NormalTok{),}
    \FunctionTok{list}\NormalTok{(}\AttributeTok{color\_fondo =} \StringTok{"\#EEEEEE"}\NormalTok{, }\AttributeTok{color\_borde =} \StringTok{"\#999999"}\NormalTok{, }\AttributeTok{alpha =} \FloatTok{0.9}\NormalTok{, }\AttributeTok{estilo =} \StringTok{"round"}\NormalTok{, }\AttributeTok{padding =} \FloatTok{0.3}\NormalTok{)}
\NormalTok{  ), }\DecValTok{1}\NormalTok{)[[}\DecValTok{1}\NormalTok{]],}

  \CommentTok{\# Aleatorizar el sector destacado (explode)}
  \AttributeTok{sector\_destacado =} \FunctionTok{sample}\NormalTok{(}\DecValTok{0}\SpecialCharTok{:}\DecValTok{4}\NormalTok{, }\DecValTok{1}\NormalTok{)}
\NormalTok{)}

\CommentTok{\# Código Python para generar el gráfico circular}
\NormalTok{codigo\_python }\OtherTok{\textless{}{-}} \FunctionTok{paste0}\NormalTok{(}\StringTok{"}
\StringTok{import matplotlib}
\StringTok{matplotlib.use(\textquotesingle{}Agg\textquotesingle{})  \# Usar backend no interactivo}
\StringTok{import matplotlib.pyplot as plt}
\StringTok{import numpy as np}

\StringTok{\# Datos para el gráfico {-} usamos los porcentajes para mostrar que se calcularon para asegurar coherencia}
\StringTok{labels = [\textquotesingle{}"}\NormalTok{, bien1, }\StringTok{" y "}\NormalTok{, bien3, }\StringTok{"\textquotesingle{}, \textquotesingle{}Solo "}\NormalTok{, bien3, }\StringTok{"\textquotesingle{}, \textquotesingle{}Solo "}\NormalTok{, bien2, }\StringTok{"\textquotesingle{},}
\StringTok{          \textquotesingle{}Solo "}\NormalTok{, bien1, }\StringTok{"\textquotesingle{}, \textquotesingle{}"}\NormalTok{, bien1, }\StringTok{" y "}\NormalTok{, bien2, }\StringTok{"\textquotesingle{}]}

\StringTok{\# Usamos porcentajes\_mostrar para asegurar coherencia visual con los cálculos          }
\StringTok{sizes = ["}\NormalTok{, porcentajes\_mostrar}\SpecialCharTok{$}\NormalTok{bien1\_bien3, }\StringTok{", "}\NormalTok{, porcentajes\_mostrar}\SpecialCharTok{$}\NormalTok{solo\_bien3, }\StringTok{", "}\NormalTok{, porcentajes\_mostrar}\SpecialCharTok{$}\NormalTok{solo\_bien2, }\StringTok{",}
\StringTok{         "}\NormalTok{, porcentajes\_mostrar}\SpecialCharTok{$}\NormalTok{solo\_bien1, }\StringTok{", "}\NormalTok{, porcentajes\_mostrar}\SpecialCharTok{$}\NormalTok{bien1\_bien2, }\StringTok{"]}

\StringTok{colors = [\textquotesingle{}"}\NormalTok{, paleta\_seleccionada[}\DecValTok{1}\NormalTok{], }\StringTok{"\textquotesingle{}, \textquotesingle{}"}\NormalTok{, paleta\_seleccionada[}\DecValTok{2}\NormalTok{], }\StringTok{"\textquotesingle{},}
\StringTok{          \textquotesingle{}"}\NormalTok{, paleta\_seleccionada[}\DecValTok{3}\NormalTok{], }\StringTok{"\textquotesingle{}, \textquotesingle{}"}\NormalTok{, paleta\_seleccionada[}\DecValTok{4}\NormalTok{], }\StringTok{"\textquotesingle{},}
\StringTok{          \textquotesingle{}"}\NormalTok{, paleta\_seleccionada[}\DecValTok{5}\NormalTok{], }\StringTok{"\textquotesingle{}]}

\StringTok{\# Crear explode para destacar un sector específico}
\StringTok{explode = [0, 0, 0, 0, 0]}
\StringTok{explode["}\NormalTok{, estilo\_grafico}\SpecialCharTok{$}\NormalTok{sector\_destacado, }\StringTok{"] = 0.03  \# Destacar el sector seleccionado}

\StringTok{\# Crear figura con más espacio para acomodar las etiquetas externas}
\StringTok{plt.figure(figsize=(7, 6))  \# Aumentar tamaño para dar espacio a las etiquetas}

\StringTok{\# Crear gráfico circular con estilos personalizados}
\StringTok{wedges, texts, autotexts = plt.pie(}
\StringTok{    sizes,}
\StringTok{    explode=explode,}
\StringTok{    colors=colors,}
\StringTok{    shadow="}\NormalTok{, }\FunctionTok{ifelse}\NormalTok{(estilo\_grafico}\SpecialCharTok{$}\NormalTok{sombra, }\StringTok{"True"}\NormalTok{, }\StringTok{"False"}\NormalTok{), }\StringTok{",}
\StringTok{    startangle="}\NormalTok{, estilo\_grafico}\SpecialCharTok{$}\NormalTok{angulo\_inicio, }\StringTok{",  \# Ángulo de inicio aleatorizado}
\StringTok{    autopct=\textquotesingle{}\%d\%\%\textquotesingle{},  \# Siempre mostrar porcentajes enteros para evitar confusiones}
\StringTok{    pctdistance=0.75,  \# Ubicar porcentajes más cerca del borde exterior}
\StringTok{    wedgeprops=\{\textquotesingle{}edgecolor\textquotesingle{}: \textquotesingle{}"}\NormalTok{, estilo\_grafico}\SpecialCharTok{$}\NormalTok{borde}\SpecialCharTok{$}\NormalTok{color, }\StringTok{"\textquotesingle{}, \textquotesingle{}linewidth\textquotesingle{}: "}\NormalTok{, estilo\_grafico}\SpecialCharTok{$}\NormalTok{borde}\SpecialCharTok{$}\NormalTok{grosor, }\StringTok{"\},}
\StringTok{    textprops=\{\textquotesingle{}fontsize\textquotesingle{}: 10, \textquotesingle{}fontweight\textquotesingle{}: \textquotesingle{}bold\textquotesingle{}, \textquotesingle{}color\textquotesingle{}: \textquotesingle{}white\textquotesingle{}\}}
\StringTok{)}

\StringTok{\# Configuración estética de los textos de porcentaje con recuadros}
\StringTok{for autotext in autotexts:}
\StringTok{    autotext.set\_fontsize(9)  \# Tamaño de fuente para porcentajes}
\StringTok{    autotext.set\_weight(\textquotesingle{}bold\textquotesingle{})}
\StringTok{    \# Crear un recuadro semitransparente para los porcentajes}
\StringTok{"}\NormalTok{, }\ControlFlowTok{if}\NormalTok{(estilo\_grafico}\SpecialCharTok{$}\NormalTok{estilo\_porcentaje}\SpecialCharTok{$}\NormalTok{color\_fondo }\SpecialCharTok{!=} \StringTok{"none"}\NormalTok{) }\FunctionTok{paste0}\NormalTok{(}\StringTok{"    bbox\_props = \{\textquotesingle{}boxstyle\textquotesingle{}: \textquotesingle{}"}\NormalTok{, estilo\_grafico}\SpecialCharTok{$}\NormalTok{estilo\_porcentaje}\SpecialCharTok{$}\NormalTok{estilo, }\StringTok{"\textquotesingle{},}
\StringTok{                 \textquotesingle{}facecolor\textquotesingle{}: \textquotesingle{}"}\NormalTok{, estilo\_grafico}\SpecialCharTok{$}\NormalTok{estilo\_porcentaje}\SpecialCharTok{$}\NormalTok{color\_fondo, }\StringTok{"\textquotesingle{},}
\StringTok{                 \textquotesingle{}alpha\textquotesingle{}: "}\NormalTok{, estilo\_grafico}\SpecialCharTok{$}\NormalTok{estilo\_porcentaje}\SpecialCharTok{$}\NormalTok{alpha, }\StringTok{",}
\StringTok{                 \textquotesingle{}edgecolor\textquotesingle{}: \textquotesingle{}none\textquotesingle{}\}}
\StringTok{    autotext.set\_bbox(bbox\_props)"}\NormalTok{), }\StringTok{"}

\StringTok{\# Eliminar los textos del pie (los reemplazaremos con etiquetas externas)}
\StringTok{for text in texts:}
\StringTok{    text.set\_visible(False)}

\StringTok{\# Crear etiquetas externas con líneas conectoras}
\StringTok{bbox\_props = \{\textquotesingle{}boxstyle\textquotesingle{}: \textquotesingle{}"}\NormalTok{, estilo\_grafico}\SpecialCharTok{$}\NormalTok{estilo\_etiquetas}\SpecialCharTok{$}\NormalTok{estilo, }\StringTok{"\textquotesingle{},}
\StringTok{             \textquotesingle{}facecolor\textquotesingle{}: \textquotesingle{}"}\NormalTok{, estilo\_grafico}\SpecialCharTok{$}\NormalTok{estilo\_etiquetas}\SpecialCharTok{$}\NormalTok{color\_fondo, }\StringTok{"\textquotesingle{},}
\StringTok{             \textquotesingle{}edgecolor\textquotesingle{}: \textquotesingle{}"}\NormalTok{, estilo\_grafico}\SpecialCharTok{$}\NormalTok{estilo\_etiquetas}\SpecialCharTok{$}\NormalTok{color\_borde, }\StringTok{"\textquotesingle{},}
\StringTok{             \textquotesingle{}alpha\textquotesingle{}: "}\NormalTok{, estilo\_grafico}\SpecialCharTok{$}\NormalTok{estilo\_etiquetas}\SpecialCharTok{$}\NormalTok{alpha, }\StringTok{"\}}

\StringTok{\# Calcular posiciones de etiquetas externas optimizadas para evitar solapamiento}
\StringTok{def get\_label\_position(angle\_rad, wedge\_size):}
\StringTok{    \# Ajustar distancia basada en el tamaño del sector}
\StringTok{    \# Sectores más pequeños tienen etiquetas más alejadas para evitar solapamiento}
\StringTok{    distance\_factor = 1.25 if wedge\_size \textless{} 15 else 1.15}

\StringTok{    \# Determinar coordenadas}
\StringTok{    x = distance\_factor * np.cos(angle\_rad)}
\StringTok{    y = distance\_factor * np.sin(angle\_rad)}

\StringTok{    \# Ajustar alineación basada en cuadrante}
\StringTok{    if x \textless{} 0:}
\StringTok{        ha = \textquotesingle{}right\textquotesingle{}}
\StringTok{    else:}
\StringTok{        ha = \textquotesingle{}left\textquotesingle{}}

\StringTok{    if y \textless{} 0:}
\StringTok{        va = \textquotesingle{}top\textquotesingle{}}
\StringTok{    else:}
\StringTok{        va = \textquotesingle{}bottom\textquotesingle{}}

\StringTok{    \# Para sectores muy pequeños, ajustar más la posición para evitar solapamiento}
\StringTok{    if wedge\_size \textless{} 10:}
\StringTok{        x *= 1.1}
\StringTok{        y *= 1.1}

\StringTok{    return x, y, ha, va}

\StringTok{\# Variables de Python (no podemos acceder a las variables de R directamente)}
\StringTok{\# Añadimos un marcador de categoría simple}
\StringTok{labels\_with\_info = [}
\StringTok{    f\textquotesingle{}\{labels[0]\} (categoría 1)\textquotesingle{},}
\StringTok{    f\textquotesingle{}\{labels[1]\} (categoría 2)\textquotesingle{},}
\StringTok{    f\textquotesingle{}\{labels[2]\} (categoría 3)\textquotesingle{},}
\StringTok{    f\textquotesingle{}\{labels[3]\} (categoría 4)\textquotesingle{},}
\StringTok{    f\textquotesingle{}\{labels[4]\} (categoría 5)\textquotesingle{}}
\StringTok{]}

\StringTok{\# Colocar etiquetas externas con líneas conectoras}
\StringTok{for i, wedge in enumerate(wedges):}
\StringTok{    \# Calcular ángulo medio del sector}
\StringTok{    ang = (wedge.theta1 + wedge.theta2) / 2}
\StringTok{    ang\_rad = np.radians(ang)}

\StringTok{    \# Calcular coordenadas para el inicio de la línea conectora (en el borde del sector)}
\StringTok{    \# El factor 0.85 asegura que la línea empiece cerca del borde del sector}
\StringTok{    x\_edge = 0.85 * np.cos(ang\_rad)}
\StringTok{    y\_edge = 0.85 * np.sin(ang\_rad)}

\StringTok{    \# Obtener posición optimizada para la etiqueta}
\StringTok{    x\_label, y\_label, ha, va = get\_label\_position(ang\_rad, sizes[i])}

\StringTok{    \# Dibujar línea conectora}
\StringTok{    con = plt.annotate(\textquotesingle{}\textquotesingle{},}
\StringTok{                      xy=(x\_edge, y\_edge),  \# Inicio (en el borde del sector)}
\StringTok{                      xytext=(x\_label * 0.95, y\_label * 0.95),  \# Fin (cerca de la etiqueta)}
\StringTok{                      arrowprops=dict(arrowstyle=\textquotesingle{}{-}\textquotesingle{}, color=\textquotesingle{}"}\NormalTok{, estilo\_grafico}\SpecialCharTok{$}\NormalTok{estilo\_etiquetas}\SpecialCharTok{$}\NormalTok{color\_borde, }\StringTok{"\textquotesingle{}, lw=0.8))}

\StringTok{    \# Añadir etiqueta con recuadro personalizado {-} mostramos solo el texto básico sin el conteo}
\StringTok{    plt.text(x\_label, y\_label, labels[i],}
\StringTok{            fontsize=8,}
\StringTok{            fontweight=\textquotesingle{}bold\textquotesingle{},}
\StringTok{            ha=ha,}
\StringTok{            va=va,}
\StringTok{            bbox=bbox\_props,}
\StringTok{            zorder=10)}

\StringTok{\# Asegurar que el gráfico sea un círculo}
\StringTok{plt.axis(\textquotesingle{}equal\textquotesingle{})}

\StringTok{\# Añadir título en la parte inferior del gráfico}
\StringTok{plt.figtext(0.5, 0.01,  \# Posición x=0.5 (centro), y=0.01 (parte inferior)}
\StringTok{           \textquotesingle{}Distribución de bienes\textquotesingle{},}
\StringTok{           fontsize=12,}
\StringTok{           fontweight=\textquotesingle{}bold\textquotesingle{},}
\StringTok{           color=\textquotesingle{}\#333333\textquotesingle{},}
\StringTok{           ha=\textquotesingle{}center\textquotesingle{})  \# Alineación horizontal centrada}

\StringTok{\# Ajustar los márgenes para dejar espacio a las etiquetas externas}
\StringTok{plt.tight\_layout(pad=1.5, rect=[0, 0.05, 1, 0.95])  \# Ajustar rect para dar más espacio}

\StringTok{\# Guardar en múltiples formatos para asegurar compatibilidad}
\StringTok{plt.savefig(\textquotesingle{}grafico\_circular.png\textquotesingle{}, dpi=150, bbox\_inches=\textquotesingle{}tight\textquotesingle{},}
\StringTok{           transparent=True, format=\textquotesingle{}png\textquotesingle{})}
\StringTok{plt.savefig(\textquotesingle{}grafico\_circular.pdf\textquotesingle{}, dpi=150, bbox\_inches=\textquotesingle{}tight\textquotesingle{},}
\StringTok{           transparent=True, format=\textquotesingle{}pdf\textquotesingle{})}
\StringTok{plt.close()}
\StringTok{"}\NormalTok{)}

\CommentTok{\# Ejecutar código Python para generar la figura}
\FunctionTok{py\_run\_string}\NormalTok{(codigo\_python)}
\end{Highlighting}
\end{Shaded}

\section{Question}\label{question}

En La Tebaida se realizó un(a) sondeo a un grupo de personas de un(a)
compañía sobre el tipo de propiedades que poseen. Los(las) estadísticas
se presentan en la gráfica.

\includegraphics[width=0.7\linewidth,height=\textheight,keepaspectratio]{grafico_circular.png}

\subsection{Answerlist}\label{answerlist}

\begin{itemize}
\tightlist
\item
  Según el gráfico, si 63 personas de la compañía poseen solo automóvil,
  ¿cuántos personas poseen solo apartamento?
\item
  Si se escoge una persona del grupo al azar, ¿cuál es la probabilidad
  de que tenga solo vivienda? (Exprese la respuesta como un número
  decimal con 4 decimales, por ejemplo: 0.2500)
\item
  ¿Cuántas personas de la compañía fueron encuestadas en total?
\end{itemize}

\section{Solution}\label{solution}

Para resolver este problema, necesitamos analizar la información
proporcionada en el gráfico circular y aplicar conceptos de
proporcionalidad y probabilidad.

\subsubsection{Pregunta 1: Número de personas que poseen solo
apartamento}\label{pregunta-1-nuxfamero-de-personas-que-poseen-solo-apartamento}

Sabemos que 63 personas poseen solo automóvil, lo que representa el 9\%
del total según el gráfico.

Para calcular el total de personas:

\begin{verbatim}
Total = (valor_conocido × 100%) ÷ porcentaje_conocido
Total = (63 × 100%) ÷ 9%
Total = (63 × 100) ÷ 9
Total = 6300 ÷ 9
Total = 700.0000
Total = 700 personas
\end{verbatim}

Ahora podemos calcular cuántas personas poseen solo apartamento:

\begin{itemize}
\tightlist
\item
  Según el gráfico, las personas con solo apartamento representan el
  23\% del total.
\end{itemize}

\begin{verbatim}
Número de personas = (porcentaje_buscado × total_personas) ÷ 100%
Número de personas = (23% × 700) ÷ 100%
Número de personas = (23 × 700) ÷ 100
Número de personas = 16100 ÷ 100
Número de personas = 161.0000
Número de personas = 161 personas
\end{verbatim}

\subsubsection{Pregunta 2: Probabilidad de que una persona tenga solo
vivienda}\label{pregunta-2-probabilidad-de-que-una-persona-tenga-solo-vivienda}

La probabilidad se calcula dividiendo el número de casos favorables
entre el número total de casos posibles:

\begin{itemize}
\tightlist
\item
  Número de personas con solo vivienda = 210 personas
\item
  Total de personas = 700 personas
\item
  Probabilidad = 210 ÷ 700
\item
  Probabilidad = 0.3000
\end{itemize}

\subsubsection{Pregunta 3: Total de personas encuestadas en la
compañía}\label{pregunta-3-total-de-personas-encuestadas-en-la-compauxf1uxeda}

Como calculamos en la primera pregunta, el total de personas encuestadas
en la compañía es 700 personas.

Verificación de la coherencia de datos: * automóvil y apartamento: 161
personas (23\% del total) * Solo apartamento: 161 personas (23\% del
total) * Solo vivienda: 210 personas (30\% del total) * Solo automóvil:
63 personas (9\% del total) * automóvil y vivienda: 105 personas (15\%
del total)

La suma de todas estas categorías es 700 personas, que coincide
exactamente con nuestro total calculado de 700 personas.

\subsection{Answerlist}\label{answerlist-1}

\begin{itemize}
\tightlist
\item
  La respuesta es 161.
\item
  La respuesta es 0.3.
\item
  La respuesta es 700.
\end{itemize}

\section{Meta-information}\label{meta-information}

exname: proporciones\_diagrama\_circular extype: cloze exclozetype:
num\textbar num\textbar num exsolution: 161\textbar0.3\textbar700 extol:
0\textbar0.0001\textbar0 exsection:
Estadística\textbar Proporciones\textbar Interpretación de gráficos

\end{document}
